%% Preámbulo común de las charlas (Beamer, pdflatex) · charlas-charlie, estilo v\estiloversion
%% Antes de %% Preámbulo común de las charlas (Beamer, pdflatex) · charlas-charlie, estilo v\estiloversion
%% Antes de %% Preámbulo común de las charlas (Beamer, pdflatex) · charlas-charlie, estilo v\estiloversion
%% Antes de %% Preámbulo común de las charlas (Beamer, pdflatex) · charlas-charlie, estilo v\estiloversion
%% Antes de \input{preambulo}, opcionalmente (los valores por defecto van primero):
%%   \def\tema{bloques}        bloques | metropolis | clasico | tufte | oscuro | editorial
%%   \def\idioma{english}      english | spanish (babel)
%%   \def\logos{harmonie,bsc}  lista de logos: harmonie, harmonie-favicon, bsc, upc; \def\logos{} los quita
%%   \def\logoportada{auto}    auto (color en fondo claro, blanco en oscuro) | color | azul | negro | blanco | tenue
%%   \def\logopie{auto}        igual, para el pie de las slides de contenido
%%   \def\tamlogos{normal}     pequeno | normal | grande
%%   \def\letra{palatino}      palatino | charter | sans (Helvetica) | lmodern
%%   \def\fondoportada{}       vacío (sin fondo) | pulmones | nariz | pulmones-original | malla | remolinos |
%%                             espectro | curvas | ruta a una imagen propia (p. ej. figs/fondo.png)
%%   \def\fondoseccion{...}    fondo de los separadores (por defecto, el mismo que la portada)
%%   \def\progreso{numeros}    numeros (n / N) | barra (barra fina en el pie) | ambos
%%   \def\sangrado{poco}       poco (1,2 cm) | mucho (2,6 cm) | 0cm (entero) | cualquier medida
%%   \def\connotas{}           notas del orador en pantalla aparte (apagado por defecto)
%% Desde latexmk:  latexmk -pdf -usepretex='\def\tema{oscuro}' -jobname=charla-oscuro main.tex
\def\estiloversion{1.2.0}
\providecommand\tema{bloques}
\providecommand\idioma{english}
\providecommand\logos{harmonie,bsc}
\providecommand\fondoportada{}
\providecommand\fondoseccion{\fondoportada}
\providecommand\progreso{numeros}
\providecommand\sangrado{poco}
\providecommand\logoportada{auto}
\providecommand\logopie{auto}
\providecommand\tamlogos{normal}
\providecommand\letra{palatino}
\providecommand\estilodir{estilo/}        % ruta a esta carpeta desde la charla

\usepackage[T1]{fontenc}
\usepackage[utf8]{inputenc}
\usepackage{babel}
\babelprovide[import,main]{\idioma}
\usepackage{amsmath,amssymb}
\makeatletter
\def\letra@a{palatino}\def\letra@b{charter}\def\letra@c{sans}\def\letra@d{lmodern}
\ifx\letra\letra@b \usepackage{charter}\fi  % Charter en el texto; fórmulas en Computer Modern
\ifx\letra\letra@b\else\ifx\letra\letra@c\else\ifx\letra\letra@d\else
  \usepackage[sc]{mathpazo}              % Palatino en texto y fórmulas (igual que las figuras)
  \renewcommand\sfdefault{ppl}\fi\fi\fi  % lo que pida \sffamily también sale en Palatino
\ifx\letra\letra@c \usepackage[scaled=0.92]{helvet}\renewcommand\familydefault{\sfdefault}\fi
\ifx\letra\letra@d \usepackage{lmodern}\fi

\ifx\letra\letra@c\else\usefonttheme{serif}\fi
\makeatother
\usefonttheme{professionalfonts}
\usepackage{microtype}
\usepackage{booktabs,multirow,array}
\usepackage{graphicx}
\usepackage{tikz}
\usetikzlibrary{positioning,arrows.meta,calc,fit,backgrounds,matrix,fadings,patterns,decorations.pathreplacing}

\setbeamertemplate{navigation symbols}{}
\ifdefined\connotas \setbeameroption{show notes on second screen=right}\fi

\makeatletter
%% ---- Fondos ---------------------------------------------------------------
%% \fondo[alto]{nombre}: imagen de fondo con opacidad \opacidadfondo, a la derecha. Lo usan la portada
%% (\fondoportada) y los separadores (\fondoseccion). \def\sinfondo{} lo quita todo.
%%   - <nombre>.tex en estilo/ (pulmones, nariz, pulmones-original): low-poly TikZ teñido con `motivo`,
%%     grande y sangrando por el borde derecho;
%%   - motivos/motivo-<nombre>.pdf: en la esquina inferior derecha;
%%   - cualquier otra ruta que exista: la imagen, abajo a la derecha.
\newsavebox\fondo@tmp
\def\sangrado@poco{1.2cm}\def\sangrado@mucho{2.6cm}
\newcommand\fondo[2][0.9]{\ifdefined\sinfondo\else
  \@ifundefined{sangrado@\sangrado}{\edef\sangrado@d{\sangrado}}{\edef\sangrado@d{\csname sangrado@\sangrado\endcsname}}\edef\fondo@n{#2}\ifx\fondo@n\@empty\else
  \IfFileExists{\estilodir\fondo@n.tex}{%
    \@ifundefined{fondo@box@\fondo@n}{%
      \expandafter\newsavebox\csname fondo@box@\fondo@n\endcsname
      \global\expandafter\sbox\csname fondo@box@\fondo@n\endcsname{\input{\estilodir\fondo@n.tex}}}{}%
    \sbox\fondo@tmp{\resizebox{!}{#1\paperheight}{\expandafter\usebox\csname fondo@box@\fondo@n\endcsname}}%
    \def\fondo@ancla{anchor=east,inner sep=0pt] at ([xshift=\sangrado@d,yshift=-0.3cm]current page.east}%
    \ifdim\wd\fondo@tmp>0.62\paperwidth % apaisado (nariz): más estrecho y abajo, lejos del título
      \sbox\fondo@tmp{\resizebox{0.56\paperwidth}{!}{\expandafter\usebox\csname fondo@box@\fondo@n\endcsname}}%
      \def\fondo@ancla{anchor=south east,inner sep=0pt] at ([xshift=\sangrado@d,yshift=0.2cm]current page.south east}\fi}{%
  \IfFileExists{\estilodir motivos/motivo-\fondo@n.pdf}{%
    \sbox\fondo@tmp{\includegraphics[height=0.8\paperheight,width=0.4\paperwidth,keepaspectratio]%
      {\estilodir motivos/motivo-\fondo@n.pdf}}%
    \def\fondo@ancla{anchor=south east,inner sep=0pt] at ([xshift=-0.3cm,yshift=0.5cm]current page.south east}}{%
  \IfFileExists{\fondo@n}{%
    \sbox\fondo@tmp{\includegraphics[height=#1\paperheight,width=0.55\paperwidth,keepaspectratio]{\fondo@n}}%
    \def\fondo@ancla{anchor=south east,inner sep=0pt] at ([xshift=-0.3cm,yshift=0.4cm]current page.south east}}{%
  \PackageWarning{charlas-charlie}{Fondo '\fondo@n' no encontrado}\sbox\fondo@tmp{}%
    \def\fondo@ancla{inner sep=0pt] at (current page.center}}}}%
  % transparency group dentro de un \tikz interno: suelto en un overlay no se dibuja
  \begin{tikzpicture}[remember picture,overlay]
    \expandafter\node\expandafter[\fondo@ancla)
      {\tikz{\begin{scope}[transparency group,opacity=\opacidadfondo]
        \node[inner sep=0pt]{\usebox\fondo@tmp};\end{scope}}};
  \end{tikzpicture}\fi\fi}

%% Adorno pequeño en la esquina superior izquierda de la portada (\def\adornoportada{remolinos}
%% o malla, espectro, curvas). Vacío = sin adorno. Los motivos PDF son blancos: solo para portadas oscuras.
\providecommand\adornoportada{}
\newcommand\adorno{\ifdefined\sinfondo\else\ifx\adornoportada\@empty\else
  \begin{tikzpicture}[remember picture,overlay]
    \node[anchor=north west,inner sep=0pt] at ([xshift=-0.5cm,yshift=0.6cm]current page.north west)
      {\tikz{\begin{scope}[transparency group,opacity=\opacidadfondo]
        \node[inner sep=0pt]{\includegraphics[height=0.42\paperheight]{\estilodir motivos/motivo-\adornoportada.pdf}};
      \end{scope}}};
  \end{tikzpicture}\fi\fi}

%% ---- Logos ----------------------------------------------------------------
%% Ficheros: estilo/logos/logo-<nombre>-<variante>.png (variantes color, azul, negro, blanco).
%% \filalogos{alto}{variante}{separación}{auto}: los logos de \logos en fila. 'auto' toma la variante que
%% fija el tema (\autoportada, \autopie: color en fondo claro, blanco en oscuro); 'tenue' = auto al 45 %.
\def\logo@auto{auto}\def\logo@tenue{tenue}
% compensación óptica: el BSC tiene cuatro líneas y el favicon es casi cuadrado
\def\escalalogo@bsc{1.3}\def\escalalogo@upc{1.15}\@namedef{escalalogo@harmonie-favicon}{1.4}
\def\tamlogos@pequeno{{0.28cm}{0.6cm}}\def\tamlogos@normal{{0.35cm}{0.75cm}}\def\tamlogos@grande{{0.45cm}{0.95cm}}
\def\tam@pie{\expandafter\expandafter\expandafter\@firstoftwo\csname tamlogos@\tamlogos\endcsname}
\def\tam@portada{\expandafter\expandafter\expandafter\@secondoftwo\csname tamlogos@\tamlogos\endcsname}
\newcommand\filalogos[4]{\begingroup\def\logo@sep{}\edef\logo@v{#2}\def\logo@op{1}%
  \ifx\logo@v\logo@tenue\edef\logo@v{#4}\def\logo@op{0.45}\fi
  \ifx\logo@v\logo@auto\edef\logo@v{#4}\fi
  \@for\logo@n:=\logos\do{\ifx\logo@n\@empty\else
    \@ifundefined{escalalogo@\logo@n}{\def\logo@k{1}}{\edef\logo@k{\csname escalalogo@\logo@n\endcsname}}%
    \logo@sep\tikz[baseline=0pt]\node[opacity=\logo@op,inner sep=0pt,anchor=south west]
      {\includegraphics[height=\logo@k\dimexpr#1\relax]{\estilodir logos/logo-\logo@n-\logo@v.png}};%
    \def\logo@sep{\hspace{#3}}\fi}\endgroup}
\providecommand\autoportada{color}\providecommand\autopie{color}
%% En la portada: grandes, abajo a la izquierda, alineados con el texto
\newcommand\logosportada{\ifx\logos\@empty\else
  \begin{tikzpicture}[remember picture,overlay]
    \node[anchor=south west,inner sep=0pt] at ([xshift=\margenizq,yshift=0.6cm]current page.south west)
      {\filalogos{\tam@portada}{\logoportada}{0.7cm}{\autoportada}};
  \end{tikzpicture}\fi}
%% Al pie de la portada: hueco a la medida de los logos (el bloque de título queda centrado en lo que sobra)
\newcommand\reservalogos{\ifx\logos\@empty\else\vspace*{\dimexpr1.3\dimexpr\tam@portada\relax+0.4cm\relax}\fi}
%% En el pie de las slides de contenido: pequeños. No van en separadores ni en slides-frase.
\newif\ifsinlogos
\BeforeBeginEnvironment{frame}{\global\sinlogosfalse}
\newcommand\ocultarlogos{\global\sinlogostrue}
\newcommand\logospie{\ifsinlogos\else\ifx\logos\@empty\else
  \raisebox{-0.1cm}{\filalogos{\tam@pie}{\logopie}{0.3cm}{\autopie}}\hspace{0.45cm}\fi\fi}

%% ---- Progreso -------------------------------------------------------------
%% \barraprogreso: línea fina a todo el ancho, llena hasta n/N, en \colorprogreso (lo puede fijar el tema)
\providecommand\colorprogreso{acento}
\def\prog@numeros{numeros}\def\prog@barra{barra}
\newcommand\barraprogreso{%
  \pgfmathsetmacro\prog@f{min(1,\insertframenumber/max(1,\inserttotalframenumber))}%
  \begin{tikzpicture}
    \fill[regla] (0,0) rectangle (\paperwidth,1.2pt);
    \fill[\colorprogreso] (0,0) rectangle (\prog@f\paperwidth,1.2pt);
  \end{tikzpicture}}
\newcommand\numeropie{\ifx\progreso\prog@barra\else\insertframenumber\,/\,\inserttotalframenumber\fi}
\newcommand\barrapie{\ifx\progreso\prog@numeros\else\par\nointerlineskip\barraprogreso\fi}
\makeatother

%% Paleta (cada tema define estos colores; mismos HEX que charla-<tema>.mplstyle):
%%   fondo, tinta (texto), tenue (texto secundario, ejes), regla (líneas finas),
%%   acento (énfasis), textochip (texto sobre relleno de serie), motivo (color de los fondos low-poly), serieA..serieD (casos en figuras)
\colorlet{textochip}{white} % texto sobre un relleno de serie (el tema oscuro lo cambia a fondo)
\input{tema-\tema.tex}
\providecommand\opacidadfondo{0.22}
% colores de pulmones-original.tex (la versión con halos y niebla)
\colorlet{lpedge}{motivo}\colorlet{lpfill}{motivo}\colorlet{lpglow}{motivo}
\colorlet{lpfloat}{motivo}\colorlet{lpfog}{motivo!85!black}\colorlet{lpnode}{motivo!60!black}
\colorlet{lpbg}{fondo}

%% ---- Macros comunes ------------------------------------------------------
%% Nombre canónico de un caso (versalitas; si el proyecto tiene GLOSSARY.md, de ahí): \caso{imp oss}
\newcommand\caso[1]{{\mdseries\scshape\lowercase{#1}}}  % mdseries: Palatino no tiene versalitas negritas
%% Slide-frase: el mensaje de un bloque, grande y centrado (\frase{...} dentro de un frame [c]). Sin logos.
\newcommand\frase[1]{\ocultarlogos{\centering\usebeamercolor[fg]{alerted text}\LARGE\bfseries\boldmath #1\par}}
%% Mensaje principal de una slide (una línea, en acento)
\newcommand\mensaje[1]{\par{\usebeamercolor[fg]{alerted text}\large #1\par}}
%% \figura[opciones]{fichero}{ancho}{alto}{qué va}: la figura si existe; si no, un hueco con lo que falta
\newcommand\figura[5][]{%
  \IfFileExists{#2}{\includegraphics[width=#3,height=#4,keepaspectratio,#1]{#2}}{%
    \begin{tikzpicture}
      \node[draw=tenue,dashed,rounded corners=2pt,minimum width=#3,minimum height=#4,
            inner sep=4pt,align=center,text width=#3-8pt,font=\scriptsize\color{tenue}]
           {#5\\[3pt]\texttt{\detokenize{#2}}};
    \end{tikzpicture}}}
, opcionalmente (los valores por defecto van primero):
%%   \def\tema{bloques}        bloques | metropolis | clasico | tufte | oscuro | editorial
%%   \def\idioma{english}      english | spanish (babel)
%%   \def\logos{harmonie,bsc}  lista de logos: harmonie, harmonie-favicon, bsc, upc; \def\logos{} los quita
%%   \def\logoportada{auto}    auto (color en fondo claro, blanco en oscuro) | color | azul | negro | blanco | tenue
%%   \def\logopie{auto}        igual, para el pie de las slides de contenido
%%   \def\tamlogos{normal}     pequeno | normal | grande
%%   \def\letra{palatino}      palatino | charter | sans (Helvetica) | lmodern
%%   \def\fondoportada{}       vacío (sin fondo) | pulmones | nariz | pulmones-original | malla | remolinos |
%%                             espectro | curvas | ruta a una imagen propia (p. ej. figs/fondo.png)
%%   \def\fondoseccion{...}    fondo de los separadores (por defecto, el mismo que la portada)
%%   \def\progreso{numeros}    numeros (n / N) | barra (barra fina en el pie) | ambos
%%   \def\sangrado{poco}       poco (1,2 cm) | mucho (2,6 cm) | 0cm (entero) | cualquier medida
%%   \def\connotas{}           notas del orador en pantalla aparte (apagado por defecto)
%% Desde latexmk:  latexmk -pdf -usepretex='\def\tema{oscuro}' -jobname=charla-oscuro main.tex
\def\estiloversion{1.2.0}
\providecommand\tema{bloques}
\providecommand\idioma{english}
\providecommand\logos{harmonie,bsc}
\providecommand\fondoportada{}
\providecommand\fondoseccion{\fondoportada}
\providecommand\progreso{numeros}
\providecommand\sangrado{poco}
\providecommand\logoportada{auto}
\providecommand\logopie{auto}
\providecommand\tamlogos{normal}
\providecommand\letra{palatino}
\providecommand\estilodir{estilo/}        % ruta a esta carpeta desde la charla

\usepackage[T1]{fontenc}
\usepackage[utf8]{inputenc}
\usepackage{babel}
\babelprovide[import,main]{\idioma}
\usepackage{amsmath,amssymb}
\makeatletter
\def\letra@a{palatino}\def\letra@b{charter}\def\letra@c{sans}\def\letra@d{lmodern}
\ifx\letra\letra@b \usepackage{charter}\fi  % Charter en el texto; fórmulas en Computer Modern
\ifx\letra\letra@b\else\ifx\letra\letra@c\else\ifx\letra\letra@d\else
  \usepackage[sc]{mathpazo}              % Palatino en texto y fórmulas (igual que las figuras)
  \renewcommand\sfdefault{ppl}\fi\fi\fi  % lo que pida \sffamily también sale en Palatino
\ifx\letra\letra@c \usepackage[scaled=0.92]{helvet}\renewcommand\familydefault{\sfdefault}\fi
\ifx\letra\letra@d \usepackage{lmodern}\fi

\ifx\letra\letra@c\else\usefonttheme{serif}\fi
\makeatother
\usefonttheme{professionalfonts}
\usepackage{microtype}
\usepackage{booktabs,multirow,array}
\usepackage{graphicx}
\usepackage{tikz}
\usetikzlibrary{positioning,arrows.meta,calc,fit,backgrounds,matrix,fadings,patterns,decorations.pathreplacing}

\setbeamertemplate{navigation symbols}{}
\ifdefined\connotas \setbeameroption{show notes on second screen=right}\fi

\makeatletter
%% ---- Fondos ---------------------------------------------------------------
%% \fondo[alto]{nombre}: imagen de fondo con opacidad \opacidadfondo, a la derecha. Lo usan la portada
%% (\fondoportada) y los separadores (\fondoseccion). \def\sinfondo{} lo quita todo.
%%   - <nombre>.tex en estilo/ (pulmones, nariz, pulmones-original): low-poly TikZ teñido con `motivo`,
%%     grande y sangrando por el borde derecho;
%%   - motivos/motivo-<nombre>.pdf: en la esquina inferior derecha;
%%   - cualquier otra ruta que exista: la imagen, abajo a la derecha.
\newsavebox\fondo@tmp
\def\sangrado@poco{1.2cm}\def\sangrado@mucho{2.6cm}
\newcommand\fondo[2][0.9]{\ifdefined\sinfondo\else
  \@ifundefined{sangrado@\sangrado}{\edef\sangrado@d{\sangrado}}{\edef\sangrado@d{\csname sangrado@\sangrado\endcsname}}\edef\fondo@n{#2}\ifx\fondo@n\@empty\else
  \IfFileExists{\estilodir\fondo@n.tex}{%
    \@ifundefined{fondo@box@\fondo@n}{%
      \expandafter\newsavebox\csname fondo@box@\fondo@n\endcsname
      \global\expandafter\sbox\csname fondo@box@\fondo@n\endcsname{\input{\estilodir\fondo@n.tex}}}{}%
    \sbox\fondo@tmp{\resizebox{!}{#1\paperheight}{\expandafter\usebox\csname fondo@box@\fondo@n\endcsname}}%
    \def\fondo@ancla{anchor=east,inner sep=0pt] at ([xshift=\sangrado@d,yshift=-0.3cm]current page.east}%
    \ifdim\wd\fondo@tmp>0.62\paperwidth % apaisado (nariz): más estrecho y abajo, lejos del título
      \sbox\fondo@tmp{\resizebox{0.56\paperwidth}{!}{\expandafter\usebox\csname fondo@box@\fondo@n\endcsname}}%
      \def\fondo@ancla{anchor=south east,inner sep=0pt] at ([xshift=\sangrado@d,yshift=0.2cm]current page.south east}\fi}{%
  \IfFileExists{\estilodir motivos/motivo-\fondo@n.pdf}{%
    \sbox\fondo@tmp{\includegraphics[height=0.8\paperheight,width=0.4\paperwidth,keepaspectratio]%
      {\estilodir motivos/motivo-\fondo@n.pdf}}%
    \def\fondo@ancla{anchor=south east,inner sep=0pt] at ([xshift=-0.3cm,yshift=0.5cm]current page.south east}}{%
  \IfFileExists{\fondo@n}{%
    \sbox\fondo@tmp{\includegraphics[height=#1\paperheight,width=0.55\paperwidth,keepaspectratio]{\fondo@n}}%
    \def\fondo@ancla{anchor=south east,inner sep=0pt] at ([xshift=-0.3cm,yshift=0.4cm]current page.south east}}{%
  \PackageWarning{charlas-charlie}{Fondo '\fondo@n' no encontrado}\sbox\fondo@tmp{}%
    \def\fondo@ancla{inner sep=0pt] at (current page.center}}}}%
  % transparency group dentro de un \tikz interno: suelto en un overlay no se dibuja
  \begin{tikzpicture}[remember picture,overlay]
    \expandafter\node\expandafter[\fondo@ancla)
      {\tikz{\begin{scope}[transparency group,opacity=\opacidadfondo]
        \node[inner sep=0pt]{\usebox\fondo@tmp};\end{scope}}};
  \end{tikzpicture}\fi\fi}

%% Adorno pequeño en la esquina superior izquierda de la portada (\def\adornoportada{remolinos}
%% o malla, espectro, curvas). Vacío = sin adorno. Los motivos PDF son blancos: solo para portadas oscuras.
\providecommand\adornoportada{}
\newcommand\adorno{\ifdefined\sinfondo\else\ifx\adornoportada\@empty\else
  \begin{tikzpicture}[remember picture,overlay]
    \node[anchor=north west,inner sep=0pt] at ([xshift=-0.5cm,yshift=0.6cm]current page.north west)
      {\tikz{\begin{scope}[transparency group,opacity=\opacidadfondo]
        \node[inner sep=0pt]{\includegraphics[height=0.42\paperheight]{\estilodir motivos/motivo-\adornoportada.pdf}};
      \end{scope}}};
  \end{tikzpicture}\fi\fi}

%% ---- Logos ----------------------------------------------------------------
%% Ficheros: estilo/logos/logo-<nombre>-<variante>.png (variantes color, azul, negro, blanco).
%% \filalogos{alto}{variante}{separación}{auto}: los logos de \logos en fila. 'auto' toma la variante que
%% fija el tema (\autoportada, \autopie: color en fondo claro, blanco en oscuro); 'tenue' = auto al 45 %.
\def\logo@auto{auto}\def\logo@tenue{tenue}
% compensación óptica: el BSC tiene cuatro líneas y el favicon es casi cuadrado
\def\escalalogo@bsc{1.3}\def\escalalogo@upc{1.15}\@namedef{escalalogo@harmonie-favicon}{1.4}
\def\tamlogos@pequeno{{0.28cm}{0.6cm}}\def\tamlogos@normal{{0.35cm}{0.75cm}}\def\tamlogos@grande{{0.45cm}{0.95cm}}
\def\tam@pie{\expandafter\expandafter\expandafter\@firstoftwo\csname tamlogos@\tamlogos\endcsname}
\def\tam@portada{\expandafter\expandafter\expandafter\@secondoftwo\csname tamlogos@\tamlogos\endcsname}
\newcommand\filalogos[4]{\begingroup\def\logo@sep{}\edef\logo@v{#2}\def\logo@op{1}%
  \ifx\logo@v\logo@tenue\edef\logo@v{#4}\def\logo@op{0.45}\fi
  \ifx\logo@v\logo@auto\edef\logo@v{#4}\fi
  \@for\logo@n:=\logos\do{\ifx\logo@n\@empty\else
    \@ifundefined{escalalogo@\logo@n}{\def\logo@k{1}}{\edef\logo@k{\csname escalalogo@\logo@n\endcsname}}%
    \logo@sep\tikz[baseline=0pt]\node[opacity=\logo@op,inner sep=0pt,anchor=south west]
      {\includegraphics[height=\logo@k\dimexpr#1\relax]{\estilodir logos/logo-\logo@n-\logo@v.png}};%
    \def\logo@sep{\hspace{#3}}\fi}\endgroup}
\providecommand\autoportada{color}\providecommand\autopie{color}
%% En la portada: grandes, abajo a la izquierda, alineados con el texto
\newcommand\logosportada{\ifx\logos\@empty\else
  \begin{tikzpicture}[remember picture,overlay]
    \node[anchor=south west,inner sep=0pt] at ([xshift=\margenizq,yshift=0.6cm]current page.south west)
      {\filalogos{\tam@portada}{\logoportada}{0.7cm}{\autoportada}};
  \end{tikzpicture}\fi}
%% Al pie de la portada: hueco a la medida de los logos (el bloque de título queda centrado en lo que sobra)
\newcommand\reservalogos{\ifx\logos\@empty\else\vspace*{\dimexpr1.3\dimexpr\tam@portada\relax+0.4cm\relax}\fi}
%% En el pie de las slides de contenido: pequeños. No van en separadores ni en slides-frase.
\newif\ifsinlogos
\BeforeBeginEnvironment{frame}{\global\sinlogosfalse}
\newcommand\ocultarlogos{\global\sinlogostrue}
\newcommand\logospie{\ifsinlogos\else\ifx\logos\@empty\else
  \raisebox{-0.1cm}{\filalogos{\tam@pie}{\logopie}{0.3cm}{\autopie}}\hspace{0.45cm}\fi\fi}

%% ---- Progreso -------------------------------------------------------------
%% \barraprogreso: línea fina a todo el ancho, llena hasta n/N, en \colorprogreso (lo puede fijar el tema)
\providecommand\colorprogreso{acento}
\def\prog@numeros{numeros}\def\prog@barra{barra}
\newcommand\barraprogreso{%
  \pgfmathsetmacro\prog@f{min(1,\insertframenumber/max(1,\inserttotalframenumber))}%
  \begin{tikzpicture}
    \fill[regla] (0,0) rectangle (\paperwidth,1.2pt);
    \fill[\colorprogreso] (0,0) rectangle (\prog@f\paperwidth,1.2pt);
  \end{tikzpicture}}
\newcommand\numeropie{\ifx\progreso\prog@barra\else\insertframenumber\,/\,\inserttotalframenumber\fi}
\newcommand\barrapie{\ifx\progreso\prog@numeros\else\par\nointerlineskip\barraprogreso\fi}
\makeatother

%% Paleta (cada tema define estos colores; mismos HEX que charla-<tema>.mplstyle):
%%   fondo, tinta (texto), tenue (texto secundario, ejes), regla (líneas finas),
%%   acento (énfasis), textochip (texto sobre relleno de serie), motivo (color de los fondos low-poly), serieA..serieD (casos en figuras)
\colorlet{textochip}{white} % texto sobre un relleno de serie (el tema oscuro lo cambia a fondo)
\input{tema-\tema.tex}
\providecommand\opacidadfondo{0.22}
% colores de pulmones-original.tex (la versión con halos y niebla)
\colorlet{lpedge}{motivo}\colorlet{lpfill}{motivo}\colorlet{lpglow}{motivo}
\colorlet{lpfloat}{motivo}\colorlet{lpfog}{motivo!85!black}\colorlet{lpnode}{motivo!60!black}
\colorlet{lpbg}{fondo}

%% ---- Macros comunes ------------------------------------------------------
%% Nombre canónico de un caso (versalitas; si el proyecto tiene GLOSSARY.md, de ahí): \caso{imp oss}
\newcommand\caso[1]{{\mdseries\scshape\lowercase{#1}}}  % mdseries: Palatino no tiene versalitas negritas
%% Slide-frase: el mensaje de un bloque, grande y centrado (\frase{...} dentro de un frame [c]). Sin logos.
\newcommand\frase[1]{\ocultarlogos{\centering\usebeamercolor[fg]{alerted text}\LARGE\bfseries\boldmath #1\par}}
%% Mensaje principal de una slide (una línea, en acento)
\newcommand\mensaje[1]{\par{\usebeamercolor[fg]{alerted text}\large #1\par}}
%% \figura[opciones]{fichero}{ancho}{alto}{qué va}: la figura si existe; si no, un hueco con lo que falta
\newcommand\figura[5][]{%
  \IfFileExists{#2}{\includegraphics[width=#3,height=#4,keepaspectratio,#1]{#2}}{%
    \begin{tikzpicture}
      \node[draw=tenue,dashed,rounded corners=2pt,minimum width=#3,minimum height=#4,
            inner sep=4pt,align=center,text width=#3-8pt,font=\scriptsize\color{tenue}]
           {#5\\[3pt]\texttt{\detokenize{#2}}};
    \end{tikzpicture}}}
, opcionalmente (los valores por defecto van primero):
%%   \def\tema{bloques}        bloques | metropolis | clasico | tufte | oscuro | editorial
%%   \def\idioma{english}      english | spanish (babel)
%%   \def\logos{harmonie,bsc}  lista de logos: harmonie, harmonie-favicon, bsc, upc; \def\logos{} los quita
%%   \def\logoportada{auto}    auto (color en fondo claro, blanco en oscuro) | color | azul | negro | blanco | tenue
%%   \def\logopie{auto}        igual, para el pie de las slides de contenido
%%   \def\tamlogos{normal}     pequeno | normal | grande
%%   \def\letra{palatino}      palatino | charter | sans (Helvetica) | lmodern
%%   \def\fondoportada{}       vacío (sin fondo) | pulmones | nariz | pulmones-original | malla | remolinos |
%%                             espectro | curvas | ruta a una imagen propia (p. ej. figs/fondo.png)
%%   \def\fondoseccion{...}    fondo de los separadores (por defecto, el mismo que la portada)
%%   \def\progreso{numeros}    numeros (n / N) | barra (barra fina en el pie) | ambos
%%   \def\sangrado{poco}       poco (1,2 cm) | mucho (2,6 cm) | 0cm (entero) | cualquier medida
%%   \def\connotas{}           notas del orador en pantalla aparte (apagado por defecto)
%% Desde latexmk:  latexmk -pdf -usepretex='\def\tema{oscuro}' -jobname=charla-oscuro main.tex
\def\estiloversion{1.2.0}
\providecommand\tema{bloques}
\providecommand\idioma{english}
\providecommand\logos{harmonie,bsc}
\providecommand\fondoportada{}
\providecommand\fondoseccion{\fondoportada}
\providecommand\progreso{numeros}
\providecommand\sangrado{poco}
\providecommand\logoportada{auto}
\providecommand\logopie{auto}
\providecommand\tamlogos{normal}
\providecommand\letra{palatino}
\providecommand\estilodir{estilo/}        % ruta a esta carpeta desde la charla

\usepackage[T1]{fontenc}
\usepackage[utf8]{inputenc}
\usepackage{babel}
\babelprovide[import,main]{\idioma}
\usepackage{amsmath,amssymb}
\makeatletter
\def\letra@a{palatino}\def\letra@b{charter}\def\letra@c{sans}\def\letra@d{lmodern}
\ifx\letra\letra@b \usepackage{charter}\fi  % Charter en el texto; fórmulas en Computer Modern
\ifx\letra\letra@b\else\ifx\letra\letra@c\else\ifx\letra\letra@d\else
  \usepackage[sc]{mathpazo}              % Palatino en texto y fórmulas (igual que las figuras)
  \renewcommand\sfdefault{ppl}\fi\fi\fi  % lo que pida \sffamily también sale en Palatino
\ifx\letra\letra@c \usepackage[scaled=0.92]{helvet}\renewcommand\familydefault{\sfdefault}\fi
\ifx\letra\letra@d \usepackage{lmodern}\fi

\ifx\letra\letra@c\else\usefonttheme{serif}\fi
\makeatother
\usefonttheme{professionalfonts}
\usepackage{microtype}
\usepackage{booktabs,multirow,array}
\usepackage{graphicx}
\usepackage{tikz}
\usetikzlibrary{positioning,arrows.meta,calc,fit,backgrounds,matrix,fadings,patterns,decorations.pathreplacing}

\setbeamertemplate{navigation symbols}{}
\ifdefined\connotas \setbeameroption{show notes on second screen=right}\fi

\makeatletter
%% ---- Fondos ---------------------------------------------------------------
%% \fondo[alto]{nombre}: imagen de fondo con opacidad \opacidadfondo, a la derecha. Lo usan la portada
%% (\fondoportada) y los separadores (\fondoseccion). \def\sinfondo{} lo quita todo.
%%   - <nombre>.tex en estilo/ (pulmones, nariz, pulmones-original): low-poly TikZ teñido con `motivo`,
%%     grande y sangrando por el borde derecho;
%%   - motivos/motivo-<nombre>.pdf: en la esquina inferior derecha;
%%   - cualquier otra ruta que exista: la imagen, abajo a la derecha.
\newsavebox\fondo@tmp
\def\sangrado@poco{1.2cm}\def\sangrado@mucho{2.6cm}
\newcommand\fondo[2][0.9]{\ifdefined\sinfondo\else
  \@ifundefined{sangrado@\sangrado}{\edef\sangrado@d{\sangrado}}{\edef\sangrado@d{\csname sangrado@\sangrado\endcsname}}\edef\fondo@n{#2}\ifx\fondo@n\@empty\else
  \IfFileExists{\estilodir\fondo@n.tex}{%
    \@ifundefined{fondo@box@\fondo@n}{%
      \expandafter\newsavebox\csname fondo@box@\fondo@n\endcsname
      \global\expandafter\sbox\csname fondo@box@\fondo@n\endcsname{\input{\estilodir\fondo@n.tex}}}{}%
    \sbox\fondo@tmp{\resizebox{!}{#1\paperheight}{\expandafter\usebox\csname fondo@box@\fondo@n\endcsname}}%
    \def\fondo@ancla{anchor=east,inner sep=0pt] at ([xshift=\sangrado@d,yshift=-0.3cm]current page.east}%
    \ifdim\wd\fondo@tmp>0.62\paperwidth % apaisado (nariz): más estrecho y abajo, lejos del título
      \sbox\fondo@tmp{\resizebox{0.56\paperwidth}{!}{\expandafter\usebox\csname fondo@box@\fondo@n\endcsname}}%
      \def\fondo@ancla{anchor=south east,inner sep=0pt] at ([xshift=\sangrado@d,yshift=0.2cm]current page.south east}\fi}{%
  \IfFileExists{\estilodir motivos/motivo-\fondo@n.pdf}{%
    \sbox\fondo@tmp{\includegraphics[height=0.8\paperheight,width=0.4\paperwidth,keepaspectratio]%
      {\estilodir motivos/motivo-\fondo@n.pdf}}%
    \def\fondo@ancla{anchor=south east,inner sep=0pt] at ([xshift=-0.3cm,yshift=0.5cm]current page.south east}}{%
  \IfFileExists{\fondo@n}{%
    \sbox\fondo@tmp{\includegraphics[height=#1\paperheight,width=0.55\paperwidth,keepaspectratio]{\fondo@n}}%
    \def\fondo@ancla{anchor=south east,inner sep=0pt] at ([xshift=-0.3cm,yshift=0.4cm]current page.south east}}{%
  \PackageWarning{charlas-charlie}{Fondo '\fondo@n' no encontrado}\sbox\fondo@tmp{}%
    \def\fondo@ancla{inner sep=0pt] at (current page.center}}}}%
  % transparency group dentro de un \tikz interno: suelto en un overlay no se dibuja
  \begin{tikzpicture}[remember picture,overlay]
    \expandafter\node\expandafter[\fondo@ancla)
      {\tikz{\begin{scope}[transparency group,opacity=\opacidadfondo]
        \node[inner sep=0pt]{\usebox\fondo@tmp};\end{scope}}};
  \end{tikzpicture}\fi\fi}

%% Adorno pequeño en la esquina superior izquierda de la portada (\def\adornoportada{remolinos}
%% o malla, espectro, curvas). Vacío = sin adorno. Los motivos PDF son blancos: solo para portadas oscuras.
\providecommand\adornoportada{}
\newcommand\adorno{\ifdefined\sinfondo\else\ifx\adornoportada\@empty\else
  \begin{tikzpicture}[remember picture,overlay]
    \node[anchor=north west,inner sep=0pt] at ([xshift=-0.5cm,yshift=0.6cm]current page.north west)
      {\tikz{\begin{scope}[transparency group,opacity=\opacidadfondo]
        \node[inner sep=0pt]{\includegraphics[height=0.42\paperheight]{\estilodir motivos/motivo-\adornoportada.pdf}};
      \end{scope}}};
  \end{tikzpicture}\fi\fi}

%% ---- Logos ----------------------------------------------------------------
%% Ficheros: estilo/logos/logo-<nombre>-<variante>.png (variantes color, azul, negro, blanco).
%% \filalogos{alto}{variante}{separación}{auto}: los logos de \logos en fila. 'auto' toma la variante que
%% fija el tema (\autoportada, \autopie: color en fondo claro, blanco en oscuro); 'tenue' = auto al 45 %.
\def\logo@auto{auto}\def\logo@tenue{tenue}
% compensación óptica: el BSC tiene cuatro líneas y el favicon es casi cuadrado
\def\escalalogo@bsc{1.3}\def\escalalogo@upc{1.15}\@namedef{escalalogo@harmonie-favicon}{1.4}
\def\tamlogos@pequeno{{0.28cm}{0.6cm}}\def\tamlogos@normal{{0.35cm}{0.75cm}}\def\tamlogos@grande{{0.45cm}{0.95cm}}
\def\tam@pie{\expandafter\expandafter\expandafter\@firstoftwo\csname tamlogos@\tamlogos\endcsname}
\def\tam@portada{\expandafter\expandafter\expandafter\@secondoftwo\csname tamlogos@\tamlogos\endcsname}
\newcommand\filalogos[4]{\begingroup\def\logo@sep{}\edef\logo@v{#2}\def\logo@op{1}%
  \ifx\logo@v\logo@tenue\edef\logo@v{#4}\def\logo@op{0.45}\fi
  \ifx\logo@v\logo@auto\edef\logo@v{#4}\fi
  \@for\logo@n:=\logos\do{\ifx\logo@n\@empty\else
    \@ifundefined{escalalogo@\logo@n}{\def\logo@k{1}}{\edef\logo@k{\csname escalalogo@\logo@n\endcsname}}%
    \logo@sep\tikz[baseline=0pt]\node[opacity=\logo@op,inner sep=0pt,anchor=south west]
      {\includegraphics[height=\logo@k\dimexpr#1\relax]{\estilodir logos/logo-\logo@n-\logo@v.png}};%
    \def\logo@sep{\hspace{#3}}\fi}\endgroup}
\providecommand\autoportada{color}\providecommand\autopie{color}
%% En la portada: grandes, abajo a la izquierda, alineados con el texto
\newcommand\logosportada{\ifx\logos\@empty\else
  \begin{tikzpicture}[remember picture,overlay]
    \node[anchor=south west,inner sep=0pt] at ([xshift=\margenizq,yshift=0.6cm]current page.south west)
      {\filalogos{\tam@portada}{\logoportada}{0.7cm}{\autoportada}};
  \end{tikzpicture}\fi}
%% Al pie de la portada: hueco a la medida de los logos (el bloque de título queda centrado en lo que sobra)
\newcommand\reservalogos{\ifx\logos\@empty\else\vspace*{\dimexpr1.3\dimexpr\tam@portada\relax+0.4cm\relax}\fi}
%% En el pie de las slides de contenido: pequeños. No van en separadores ni en slides-frase.
\newif\ifsinlogos
\BeforeBeginEnvironment{frame}{\global\sinlogosfalse}
\newcommand\ocultarlogos{\global\sinlogostrue}
\newcommand\logospie{\ifsinlogos\else\ifx\logos\@empty\else
  \raisebox{-0.1cm}{\filalogos{\tam@pie}{\logopie}{0.3cm}{\autopie}}\hspace{0.45cm}\fi\fi}

%% ---- Progreso -------------------------------------------------------------
%% \barraprogreso: línea fina a todo el ancho, llena hasta n/N, en \colorprogreso (lo puede fijar el tema)
\providecommand\colorprogreso{acento}
\def\prog@numeros{numeros}\def\prog@barra{barra}
\newcommand\barraprogreso{%
  \pgfmathsetmacro\prog@f{min(1,\insertframenumber/max(1,\inserttotalframenumber))}%
  \begin{tikzpicture}
    \fill[regla] (0,0) rectangle (\paperwidth,1.2pt);
    \fill[\colorprogreso] (0,0) rectangle (\prog@f\paperwidth,1.2pt);
  \end{tikzpicture}}
\newcommand\numeropie{\ifx\progreso\prog@barra\else\insertframenumber\,/\,\inserttotalframenumber\fi}
\newcommand\barrapie{\ifx\progreso\prog@numeros\else\par\nointerlineskip\barraprogreso\fi}
\makeatother

%% Paleta (cada tema define estos colores; mismos HEX que charla-<tema>.mplstyle):
%%   fondo, tinta (texto), tenue (texto secundario, ejes), regla (líneas finas),
%%   acento (énfasis), textochip (texto sobre relleno de serie), motivo (color de los fondos low-poly), serieA..serieD (casos en figuras)
\colorlet{textochip}{white} % texto sobre un relleno de serie (el tema oscuro lo cambia a fondo)
\input{tema-\tema.tex}
\providecommand\opacidadfondo{0.22}
% colores de pulmones-original.tex (la versión con halos y niebla)
\colorlet{lpedge}{motivo}\colorlet{lpfill}{motivo}\colorlet{lpglow}{motivo}
\colorlet{lpfloat}{motivo}\colorlet{lpfog}{motivo!85!black}\colorlet{lpnode}{motivo!60!black}
\colorlet{lpbg}{fondo}

%% ---- Macros comunes ------------------------------------------------------
%% Nombre canónico de un caso (versalitas; si el proyecto tiene GLOSSARY.md, de ahí): \caso{imp oss}
\newcommand\caso[1]{{\mdseries\scshape\lowercase{#1}}}  % mdseries: Palatino no tiene versalitas negritas
%% Slide-frase: el mensaje de un bloque, grande y centrado (\frase{...} dentro de un frame [c]). Sin logos.
\newcommand\frase[1]{\ocultarlogos{\centering\usebeamercolor[fg]{alerted text}\LARGE\bfseries\boldmath #1\par}}
%% Mensaje principal de una slide (una línea, en acento)
\newcommand\mensaje[1]{\par{\usebeamercolor[fg]{alerted text}\large #1\par}}
%% \figura[opciones]{fichero}{ancho}{alto}{qué va}: la figura si existe; si no, un hueco con lo que falta
\newcommand\figura[5][]{%
  \IfFileExists{#2}{\includegraphics[width=#3,height=#4,keepaspectratio,#1]{#2}}{%
    \begin{tikzpicture}
      \node[draw=tenue,dashed,rounded corners=2pt,minimum width=#3,minimum height=#4,
            inner sep=4pt,align=center,text width=#3-8pt,font=\scriptsize\color{tenue}]
           {#5\\[3pt]\texttt{\detokenize{#2}}};
    \end{tikzpicture}}}
, opcionalmente (los valores por defecto van primero):
%%   \def\tema{bloques}        bloques | metropolis | clasico | tufte | oscuro | editorial
%%   \def\idioma{english}      english | spanish (babel)
%%   \def\logos{harmonie,bsc}  lista de logos: harmonie, harmonie-favicon, bsc, upc; \def\logos{} los quita
%%   \def\logoportada{auto}    auto (color en fondo claro, blanco en oscuro) | color | azul | negro | blanco | tenue
%%   \def\logopie{auto}        igual, para el pie de las slides de contenido
%%   \def\tamlogos{normal}     pequeno | normal | grande
%%   \def\letra{palatino}      palatino | charter | sans (Helvetica) | lmodern
%%   \def\fondoportada{}       vacío (sin fondo) | pulmones | nariz | pulmones-original | malla | remolinos |
%%                             espectro | curvas | ruta a una imagen propia (p. ej. figs/fondo.png)
%%   \def\fondoseccion{...}    fondo de los separadores (por defecto, el mismo que la portada)
%%   \def\progreso{numeros}    numeros (n / N) | barra (barra fina en el pie) | ambos
%%   \def\sangrado{poco}       poco (1,2 cm) | mucho (2,6 cm) | 0cm (entero) | cualquier medida
%%   \def\connotas{}           notas del orador en pantalla aparte (apagado por defecto)
%% Desde latexmk:  latexmk -pdf -usepretex='\def\tema{oscuro}' -jobname=charla-oscuro main.tex
\def\estiloversion{1.2.0}
\providecommand\tema{bloques}
\providecommand\idioma{english}
\providecommand\logos{harmonie,bsc}
\providecommand\fondoportada{}
\providecommand\fondoseccion{\fondoportada}
\providecommand\progreso{numeros}
\providecommand\sangrado{poco}
\providecommand\logoportada{auto}
\providecommand\logopie{auto}
\providecommand\tamlogos{normal}
\providecommand\letra{palatino}
\providecommand\estilodir{estilo/}        % ruta a esta carpeta desde la charla

\usepackage[T1]{fontenc}
\usepackage[utf8]{inputenc}
\usepackage{babel}
\babelprovide[import,main]{\idioma}
\usepackage{amsmath,amssymb}
\makeatletter
\def\letra@a{palatino}\def\letra@b{charter}\def\letra@c{sans}\def\letra@d{lmodern}
\ifx\letra\letra@b \usepackage{charter}\fi  % Charter en el texto; fórmulas en Computer Modern
\ifx\letra\letra@b\else\ifx\letra\letra@c\else\ifx\letra\letra@d\else
  \usepackage[sc]{mathpazo}              % Palatino en texto y fórmulas (igual que las figuras)
  \renewcommand\sfdefault{ppl}\fi\fi\fi  % lo que pida \sffamily también sale en Palatino
\ifx\letra\letra@c \usepackage[scaled=0.92]{helvet}\renewcommand\familydefault{\sfdefault}\fi
\ifx\letra\letra@d \usepackage{lmodern}\fi

\ifx\letra\letra@c\else\usefonttheme{serif}\fi
\makeatother
\usefonttheme{professionalfonts}
\usepackage{microtype}
\usepackage{booktabs,multirow,array}
\usepackage{graphicx}
\usepackage{tikz}
\usetikzlibrary{positioning,arrows.meta,calc,fit,backgrounds,matrix,fadings,patterns,decorations.pathreplacing}

\setbeamertemplate{navigation symbols}{}
\ifdefined\connotas \setbeameroption{show notes on second screen=right}\fi

\makeatletter
%% ---- Fondos ---------------------------------------------------------------
%% \fondo[alto]{nombre}: imagen de fondo con opacidad \opacidadfondo, a la derecha. Lo usan la portada
%% (\fondoportada) y los separadores (\fondoseccion). \def\sinfondo{} lo quita todo.
%%   - <nombre>.tex en estilo/ (pulmones, nariz, pulmones-original): low-poly TikZ teñido con `motivo`,
%%     grande y sangrando por el borde derecho;
%%   - motivos/motivo-<nombre>.pdf: en la esquina inferior derecha;
%%   - cualquier otra ruta que exista: la imagen, abajo a la derecha.
\newsavebox\fondo@tmp
\def\sangrado@poco{1.2cm}\def\sangrado@mucho{2.6cm}
\newcommand\fondo[2][0.9]{\ifdefined\sinfondo\else
  \@ifundefined{sangrado@\sangrado}{\edef\sangrado@d{\sangrado}}{\edef\sangrado@d{\csname sangrado@\sangrado\endcsname}}\edef\fondo@n{#2}\ifx\fondo@n\@empty\else
  \IfFileExists{\estilodir\fondo@n.tex}{%
    \@ifundefined{fondo@box@\fondo@n}{%
      \expandafter\newsavebox\csname fondo@box@\fondo@n\endcsname
      \global\expandafter\sbox\csname fondo@box@\fondo@n\endcsname{\input{\estilodir\fondo@n.tex}}}{}%
    \sbox\fondo@tmp{\resizebox{!}{#1\paperheight}{\expandafter\usebox\csname fondo@box@\fondo@n\endcsname}}%
    \def\fondo@ancla{anchor=east,inner sep=0pt] at ([xshift=\sangrado@d,yshift=-0.3cm]current page.east}%
    \ifdim\wd\fondo@tmp>0.62\paperwidth % apaisado (nariz): más estrecho y abajo, lejos del título
      \sbox\fondo@tmp{\resizebox{0.56\paperwidth}{!}{\expandafter\usebox\csname fondo@box@\fondo@n\endcsname}}%
      \def\fondo@ancla{anchor=south east,inner sep=0pt] at ([xshift=\sangrado@d,yshift=0.2cm]current page.south east}\fi}{%
  \IfFileExists{\estilodir motivos/motivo-\fondo@n.pdf}{%
    \sbox\fondo@tmp{\includegraphics[height=0.8\paperheight,width=0.4\paperwidth,keepaspectratio]%
      {\estilodir motivos/motivo-\fondo@n.pdf}}%
    \def\fondo@ancla{anchor=south east,inner sep=0pt] at ([xshift=-0.3cm,yshift=0.5cm]current page.south east}}{%
  \IfFileExists{\fondo@n}{%
    \sbox\fondo@tmp{\includegraphics[height=#1\paperheight,width=0.55\paperwidth,keepaspectratio]{\fondo@n}}%
    \def\fondo@ancla{anchor=south east,inner sep=0pt] at ([xshift=-0.3cm,yshift=0.4cm]current page.south east}}{%
  \PackageWarning{charlas-charlie}{Fondo '\fondo@n' no encontrado}\sbox\fondo@tmp{}%
    \def\fondo@ancla{inner sep=0pt] at (current page.center}}}}%
  % transparency group dentro de un \tikz interno: suelto en un overlay no se dibuja
  \begin{tikzpicture}[remember picture,overlay]
    \expandafter\node\expandafter[\fondo@ancla)
      {\tikz{\begin{scope}[transparency group,opacity=\opacidadfondo]
        \node[inner sep=0pt]{\usebox\fondo@tmp};\end{scope}}};
  \end{tikzpicture}\fi\fi}

%% Adorno pequeño en la esquina superior izquierda de la portada (\def\adornoportada{remolinos}
%% o malla, espectro, curvas). Vacío = sin adorno. Los motivos PDF son blancos: solo para portadas oscuras.
\providecommand\adornoportada{}
\newcommand\adorno{\ifdefined\sinfondo\else\ifx\adornoportada\@empty\else
  \begin{tikzpicture}[remember picture,overlay]
    \node[anchor=north west,inner sep=0pt] at ([xshift=-0.5cm,yshift=0.6cm]current page.north west)
      {\tikz{\begin{scope}[transparency group,opacity=\opacidadfondo]
        \node[inner sep=0pt]{\includegraphics[height=0.42\paperheight]{\estilodir motivos/motivo-\adornoportada.pdf}};
      \end{scope}}};
  \end{tikzpicture}\fi\fi}

%% ---- Logos ----------------------------------------------------------------
%% Ficheros: estilo/logos/logo-<nombre>-<variante>.png (variantes color, azul, negro, blanco).
%% \filalogos{alto}{variante}{separación}{auto}: los logos de \logos en fila. 'auto' toma la variante que
%% fija el tema (\autoportada, \autopie: color en fondo claro, blanco en oscuro); 'tenue' = auto al 45 %.
\def\logo@auto{auto}\def\logo@tenue{tenue}
% compensación óptica: el BSC tiene cuatro líneas y el favicon es casi cuadrado
\def\escalalogo@bsc{1.3}\def\escalalogo@upc{1.15}\@namedef{escalalogo@harmonie-favicon}{1.4}
\def\tamlogos@pequeno{{0.28cm}{0.6cm}}\def\tamlogos@normal{{0.35cm}{0.75cm}}\def\tamlogos@grande{{0.45cm}{0.95cm}}
\def\tam@pie{\expandafter\expandafter\expandafter\@firstoftwo\csname tamlogos@\tamlogos\endcsname}
\def\tam@portada{\expandafter\expandafter\expandafter\@secondoftwo\csname tamlogos@\tamlogos\endcsname}
\newcommand\filalogos[4]{\begingroup\def\logo@sep{}\edef\logo@v{#2}\def\logo@op{1}%
  \ifx\logo@v\logo@tenue\edef\logo@v{#4}\def\logo@op{0.45}\fi
  \ifx\logo@v\logo@auto\edef\logo@v{#4}\fi
  \@for\logo@n:=\logos\do{\ifx\logo@n\@empty\else
    \@ifundefined{escalalogo@\logo@n}{\def\logo@k{1}}{\edef\logo@k{\csname escalalogo@\logo@n\endcsname}}%
    \logo@sep\tikz[baseline=0pt]\node[opacity=\logo@op,inner sep=0pt,anchor=south west]
      {\includegraphics[height=\logo@k\dimexpr#1\relax]{\estilodir logos/logo-\logo@n-\logo@v.png}};%
    \def\logo@sep{\hspace{#3}}\fi}\endgroup}
\providecommand\autoportada{color}\providecommand\autopie{color}
%% En la portada: grandes, abajo a la izquierda, alineados con el texto
\newcommand\logosportada{\ifx\logos\@empty\else
  \begin{tikzpicture}[remember picture,overlay]
    \node[anchor=south west,inner sep=0pt] at ([xshift=\margenizq,yshift=0.6cm]current page.south west)
      {\filalogos{\tam@portada}{\logoportada}{0.7cm}{\autoportada}};
  \end{tikzpicture}\fi}
%% Al pie de la portada: hueco a la medida de los logos (el bloque de título queda centrado en lo que sobra)
\newcommand\reservalogos{\ifx\logos\@empty\else\vspace*{\dimexpr1.3\dimexpr\tam@portada\relax+0.4cm\relax}\fi}
%% En el pie de las slides de contenido: pequeños. No van en separadores ni en slides-frase.
\newif\ifsinlogos
\BeforeBeginEnvironment{frame}{\global\sinlogosfalse}
\newcommand\ocultarlogos{\global\sinlogostrue}
\newcommand\logospie{\ifsinlogos\else\ifx\logos\@empty\else
  \raisebox{-0.1cm}{\filalogos{\tam@pie}{\logopie}{0.3cm}{\autopie}}\hspace{0.45cm}\fi\fi}

%% ---- Progreso -------------------------------------------------------------
%% \barraprogreso: línea fina a todo el ancho, llena hasta n/N, en \colorprogreso (lo puede fijar el tema)
\providecommand\colorprogreso{acento}
\def\prog@numeros{numeros}\def\prog@barra{barra}
\newcommand\barraprogreso{%
  \pgfmathsetmacro\prog@f{min(1,\insertframenumber/max(1,\inserttotalframenumber))}%
  \begin{tikzpicture}
    \fill[regla] (0,0) rectangle (\paperwidth,1.2pt);
    \fill[\colorprogreso] (0,0) rectangle (\prog@f\paperwidth,1.2pt);
  \end{tikzpicture}}
\newcommand\numeropie{\ifx\progreso\prog@barra\else\insertframenumber\,/\,\inserttotalframenumber\fi}
\newcommand\barrapie{\ifx\progreso\prog@numeros\else\par\nointerlineskip\barraprogreso\fi}
\makeatother

%% Paleta (cada tema define estos colores; mismos HEX que charla-<tema>.mplstyle):
%%   fondo, tinta (texto), tenue (texto secundario, ejes), regla (líneas finas),
%%   acento (énfasis), textochip (texto sobre relleno de serie), motivo (color de los fondos low-poly), serieA..serieD (casos en figuras)
\colorlet{textochip}{white} % texto sobre un relleno de serie (el tema oscuro lo cambia a fondo)
\input{tema-\tema.tex}
\providecommand\opacidadfondo{0.22}
% colores de pulmones-original.tex (la versión con halos y niebla)
\colorlet{lpedge}{motivo}\colorlet{lpfill}{motivo}\colorlet{lpglow}{motivo}
\colorlet{lpfloat}{motivo}\colorlet{lpfog}{motivo!85!black}\colorlet{lpnode}{motivo!60!black}
\colorlet{lpbg}{fondo}

%% ---- Macros comunes ------------------------------------------------------
%% Nombre canónico de un caso (versalitas; si el proyecto tiene GLOSSARY.md, de ahí): \caso{imp oss}
\newcommand\caso[1]{{\mdseries\scshape\lowercase{#1}}}  % mdseries: Palatino no tiene versalitas negritas
%% Slide-frase: el mensaje de un bloque, grande y centrado (\frase{...} dentro de un frame [c]). Sin logos.
\newcommand\frase[1]{\ocultarlogos{\centering\usebeamercolor[fg]{alerted text}\LARGE\bfseries\boldmath #1\par}}
%% Mensaje principal de una slide (una línea, en acento)
\newcommand\mensaje[1]{\par{\usebeamercolor[fg]{alerted text}\large #1\par}}
%% \figura[opciones]{fichero}{ancho}{alto}{qué va}: la figura si existe; si no, un hueco con lo que falta
\newcommand\figura[5][]{%
  \IfFileExists{#2}{\includegraphics[width=#3,height=#4,keepaspectratio,#1]{#2}}{%
    \begin{tikzpicture}
      \node[draw=tenue,dashed,rounded corners=2pt,minimum width=#3,minimum height=#4,
            inner sep=4pt,align=center,text width=#3-8pt,font=\scriptsize\color{tenue}]
           {#5\\[3pt]\texttt{\detokenize{#2}}};
    \end{tikzpicture}}}
